% =============================================================================
% Primitive indecomposable P_0-shadow.
%
% Conventions:
% - bar A = C[z_1,z_2]/C·1, the reduced closed Hamiltonian Lie algebra.
% - mc P  = bigoplus_{a+b>0} C·rho_{a,b}, the principal-part coadjoint module.
% - Obs^cl_open,N(I) = (Omega^*_c(I) hat-tensor R_N^{GL_N}, Q),
%   R_N = C[(phi_1)^i_j,(phi_2)^i_j] tensor Lambda((psi)^i_j),
%   Q psi = [phi_1,phi_2], Q phi_i = 0,
%   Omega^*_c the compactly supported de Rham complex on I subset R.
% - Cyc(F)_{nonempty} the cyclic words on two letters of total positive degree;
%   stably (N >= word length), trace functions of cyclic words are linearly
%   independent (prop:stable-trace-invariants).
% - Single-trace = primitive in the cumulant / connected sense; multi-trace =
%   decomposable in the symmetric algebra on cyclic-word generators.
% - tr-count filtration: the increasing filtration F_k Obs by total trace
%   count, F_0 = scalars, F_1 = single-trace, F_k = at most k traces.
% =============================================================================

\subsection{\texorpdfstring{Primitive indecomposable \(P_0\)-shadow}{Primitive indecomposable P0-shadow}}
\label{ssec:chain-level-primitive}

The primitive indecomposable \(P_0\)-shadow map
$$\pi_{\mathrm{prim}}:
  \mathrm{Obs}^{\mathrm{cl}}_{\mathrm{open},N}(I)
  \longrightarrow
  \Omega^\bullet_c(I)\widehat\otimes\bar A$$
required by Problem~\ref{prob:formal-factorization-center}, item (3),
is constructed below without recourse to any heat-kernel propagator or
finite-scale regularization.  The construction relies only on the chain
structure
of the reduced BV algebra (Construction~\ref{constr:reduced-bv-algebra}),
the de Rham contraction (Lemma~\ref{lem:derivative-jets}), and the
stable trace invariant theory of
Proposition~\ref{prop:stable-trace-invariants}.  The principal
ingredient is the cumulant projection on the free symmetric algebra
generated by primitive cyclic-word trace classes.

\begin{constr}[Trace-count filtration on the open observables]\label{constr:trace-count-filtration}
  Let $\mc R_N$ be as in Construction~\ref{constr:reduced-bv-algebra}.
  In the stable connected limit $N\to\infty$,
  Proposition~\ref{prop:stable-trace-invariants} identifies
  $$\mc R_\infty^{GL_\infty,\mathrm{st}}
    \cong\mathbf S\bigl(\mathrm{Cyc}(F)_{\mathrm{nonempty}}\bigr)$$
  with $F=\C\langle x_1,x_2\rangle$ and the cyclic words being primitive
  generators (Lemma~\ref{lem:stable-marked-traces} extends this to the
  super-cyclic case with one displayed odd $\psi$).  The
  \emph{trace-count filtration} on
  $\mathrm{Obs}^{\mathrm{cl}}_{\mathrm{open},N}(I)
   =\Omega^\bullet_c(I)\widehat\otimes\mc R_N^{GL_N}$
  is the symmetric-algebra filtration
  $F_kObs$ generated by products of at most $k$ primitive trace classes
  (cyclic words):
  $$F_0\subset F_1\subset F_2\subset\cdots
    \subset\mathrm{Obs}^{\mathrm{cl}}_{\mathrm{open},N}(I).$$
  $F_0$ is the algebra of compactly supported de Rham forms on $I$ with
  scalar coefficients (the empty-trace stratum);
  $F_1\!\big/F_0$ is the \emph{primitive} or \emph{single-trace}
  component, isomorphic as a graded vector space to
  $\Omega^\bullet_c(I)\widehat\otimes\mathrm{Cyc}(F)_{\mathrm{nonempty}}$
  in the stable limit.
\end{constr}

\begin{lem}[Trace-count filtration is preserved by $Q$ and by the open
$P_0$ bracket]\label{lem:filtration-preservation}
  In the stable connected limit, the BV differential $Q$ and the open
  $P_0$ bracket $\{-,-\}_{\mathrm{open}}$ both preserve the
  trace-count filtration of
  Construction~\ref{constr:trace-count-filtration}.
\end{lem}

\begin{proof}
  \emph{$Q$ preserves $F_k$.}  By
  Construction~\ref{constr:reduced-bv-algebra},
  $Q\psi=[\phi_1,\phi_2]$, $Q\phi_i=0$.  $Q$ is an interior derivation
  on $\mc R_N$, so on a product of $k$ traces it acts by the Leibniz
  rule and produces another sum of $k$-trace products
  (each application replaces a single $\psi$ inside a trace by
  $[\phi_1,\phi_2]$ inside the same trace; the total trace count is
  preserved).  Smearing by a de Rham form on $I$ does not change the
  trace count.  Hence $Q F_k\subset F_k$.

  \emph{The open $P_0$ bracket preserves $F_k$.}  By the canonical
  equal-time bracket of
  Remark~\ref{rmk:equal-time-locality},
  $\{(\phi_1)^i{}_j(t),(\phi_2)^k{}_l(s)\}=\delta^i_l\delta^k_j\delta(t-s)$.
  On two traces $\mathrm{Tr}\,w_1$ and $\mathrm{Tr}\,w_2$ this bracket
  produces a single trace
  $\mathrm{Tr}\,w_3$ obtained by Wick-contracting one $\phi_1$ in $w_1$
  with one $\phi_2$ in $w_2$ and joining the two cycles into one.
  This is the standard merging-of-cycles formula of Procesi--Razmyslov
  invariant theory.  In particular, the bracket of two single-trace
  classes is a single-trace class, exactly Theorem~\ref{thm:bulk-boundary}.
  By the Leibniz rule for the Poisson bracket, the bracket of a
  $k_1$-trace product with a $k_2$-trace product is a sum of products
  of total trace count exactly $(k_1+k_2-1)$, that is, it stays in
  $F_{k_1+k_2-1}\subset F_{k_1+k_2}$.  Hence the bracket preserves the
  filtration.
\end{proof}

\begin{constr}[Primitive indecomposable \(P_0\)-shadow]\label{constr:chain-level-primitive-projection}
  In the stable connected limit, the symmetric-algebra structure
  $\mc R_\infty^{GL_\infty,\mathrm{st}}\cong
  \mathbf S(\mathrm{Cyc}(F)_{\mathrm{nonempty}})$ has the natural
  cumulant projection onto the linear span of generators (the
  Hopf-algebraic primitive). Tensored with the identity on
  $\Omega^\bullet_c(I)$, this gives
  $$\pi_{\mathrm{prim}}:
    \Omega^\bullet_c(I)\widehat\otimes\mathbf S(\mathrm{Cyc}(F)_{\mathrm{nonempty}})
    \twoheadrightarrow
    \Omega^\bullet_c(I)\widehat\otimes\mathrm{Cyc}(F)_{\mathrm{nonempty}}.$$
  Composing this shadow map with the identification
  $\mathrm{Cyc}(F)_{\mathrm{nonempty}}\simeq\bar A$ on the
  Hamiltonian-trace subspace
  (the abelianization of cyclic words: the trace of a cyclic word in
  the commuting variables $\phi_1,\phi_2$, taken modulo $Q$-exact
  reordering, is exactly the polynomial Hamiltonian
  $f\in\bar A$, a stable cyclic-word/polynomial dictionary recorded
  by Construction~\ref{constr:boundary-evaluation} and
  Lemma~\ref{lem:trace-bracket-descends}), gives
  $$\pi_{\mathrm{prim}}:
    \mathrm{Obs}^{\mathrm{cl}}_{\mathrm{open},N}(I)
    \twoheadrightarrow
    \Omega^\bullet_c(I)\widehat\otimes\bar A.$$
  At finite $N$ this is a degreewise stable construction, not a
  global map $\mc R_N\to\mc R_\infty^{\mathrm{st}}$.  For a fixed
  word-degree $d$ and $N\geq d$,
  Proposition~\ref{prop:stable-trace-invariants} identifies the
  degree-$d$ trace invariants with the stable cyclic-word span, and
  the above shadow map is applied on that stable range.  Outside the
  stable range the finite-$N$ trace identities are not discarded; the
  statement is made after passing to the stable connected limit.
\end{constr}

\begin{lem}[Cumulant kernel is the closed ideal of decomposables]\label{lem:cumulant-kernel}
  The kernel of $\pi_{\mathrm{prim}}$ as a chain map is the closed
  ideal $F_{\geq 2}$ generated by products of at least two primitive
  trace classes.  Equivalently,
  $\ker\pi_{\mathrm{prim}}=\Omega^\bullet_c(I)\widehat\otimes
  \mathbf S^{\geq 2}(\mathrm{Cyc}(F)_{\mathrm{nonempty}})$ on the
  multi-trace stratum, plus
  $\Omega^\bullet_c(I)\widehat\otimes\C\cdot1$ on the empty-trace
  stratum (constants).
\end{lem}

\begin{proof}
  By construction $\pi_{\mathrm{prim}}$ is the projection onto
  generators of the symmetric algebra, which is the standard cumulant
  projection.  Its kernel is the augmentation ideal squared
  $(\mathbf S^{\geq 1})^2=\mathbf S^{\geq 2}$, complemented by the
  constants $\C\cdot1\subset F_0$.  Tensoring with the de Rham
  cosheaf $\Omega^\bullet_c(I)$ preserves the kernel because the
  symmetric-algebra split is over the discrete coefficient ring on
  cyclic-word generators.
\end{proof}

\begin{thm}[Primitive indecomposable \(P_0\)-shadow is a Q-chain map]\label{thm:chain-level-primitive-projection-Q}
  $\pi_{\mathrm{prim}}$ commutes strictly with $Q$:
  $$\pi_{\mathrm{prim}}\circ Q=Q\circ\pi_{\mathrm{prim}}.$$
\end{thm}

\begin{proof}
  By Lemma~\ref{lem:filtration-preservation}, $Q$ preserves the
  trace-count filtration $F_\bullet$.  In particular $Q$ restricts to
  endomorphisms of $F_0$, $F_1$, and the multi-trace ideal
  $F_{\geq 2}$.  Hence $Q$ preserves the symmetric-algebra splitting
  $$F_0\oplus(F_1/F_0)\oplus F_{\geq 2}\,/\,F_{\geq 2}$$
  in the associated graded sense.  $\pi_{\mathrm{prim}}$ is by
  construction the projection onto $F_1/F_0$ in this splitting; the
  scalar empty-trace summand $F_0$ is killed by the tensor factor
  $\bar A=\C[z_1,z_2]/\C\cdot 1$.  Therefore
  $\pi_{\mathrm{prim}}\circ Q$ and $Q\circ\pi_{\mathrm{prim}}$ both
  agree on each summand: on $F_1/F_0$ both equal $Q$ restricted to
  the single-trace component; on $F_0$ and on $F_{\geq 2}$ both equal
  zero.
\end{proof}

\begin{rmk}[The cumulant target as the indecomposable part]
  \label{rmk:cumulant-target-indecomposable}
  The target
  $\Omega^\bullet_c(I)\widehat\otimes\bar A$ is the
  \emph{indecomposable} or \emph{linear primitive} part of the brane
  $P_0$ factorization algebra
  $A_\partial^{\mathrm{Ham}}(I)=\widehat{\mathbf S}(\mathfrak h_I)$
  of Construction~\ref{constr:interval-fact-algebras}: it is the
  generating Lie subspace
  $\mathfrak h_I=\Omega^\bullet_c(I)\widehat\otimes\bar A$.  Equivalently,
  it is the quotient
  $\widehat{\mathbf S}(\mathfrak h_I)/\widehat{\mathbf S}^{\geq 2}(\mathfrak h_I)
  \cong\mathfrak h_I$ in the augmentation-ideal grading, after
  removal of constants.  The factorization product on the cumulant
  target side is therefore the linear (Lie) bracket of $\mathfrak h_I$,
  not a commutative algebra product; the multiplication of two
  distinct primitive elements lands in $\mathbf S^{\geq 2}$, which is
  the zero space at the indecomposable level.
\end{rmk}

\begin{thm}[Factorization compatibility of the indecomposable shadow]\label{thm:chain-level-primitive-projection-factorization}
  For disjoint intervals $I_1\sqcup I_2\subset V$ in $\R$, the
  diagram
  $$\begin{array}{ccc}
    \mathrm{Obs}^{\mathrm{cl}}_{\mathrm{open},N}(I_1)\otimes
    \mathrm{Obs}^{\mathrm{cl}}_{\mathrm{open},N}(I_2)
    & \xrightarrow{\,\mu\,} &
    \mathrm{Obs}^{\mathrm{cl}}_{\mathrm{open},N}(V) \\
    \mathrm{aug}_{\oplus}\circ
    (\pi_{\mathrm{prim}}\otimes\pi_{\mathrm{prim}})\;\downarrow & &
    \downarrow\;\pi_{\mathrm{prim}} \\
    \bigl(\Omega^\bullet_c(I_1)\widehat\otimes\bar A\bigr)\oplus
    \bigl(\Omega^\bullet_c(I_2)\widehat\otimes\bar A\bigr)
    & \xrightarrow{\,\mu_{\bar A}\,} &
    \Omega^\bullet_c(V)\widehat\otimes\bar A
  \end{array}$$
  commutes strictly modulo a homotopy $h_{\mathrm{fact}}$ that
  vanishes identically: $h_{\mathrm{fact}}\equiv 0$.

  Here $\mu$ is the brane factorization product (multiplication in
  $\mc R_N^{GL_N}$ after extension by zero from $\Omega^\bullet_c(I_j)$ to
  $\Omega^\bullet_c(V)$), the bottom horizontal $\mu_{\bar A}$ is the
  cosheaf direct-sum extension by zero
  $\mathfrak h_{I_1}\oplus\mathfrak h_{I_2}\hookrightarrow\mathfrak h_V$
  on the indecomposable (linear) part of the brane $P_0$ factorization
  algebra (cf.~Remark~\ref{rmk:cumulant-target-indecomposable}), and
  the left vertical map factors through the augmentation
  $$\bigl(\Omega^\bullet_c(I_1)\widehat\otimes\bar A\bigr)\otimes
    \bigl(\Omega^\bullet_c(I_2)\widehat\otimes\bar A\bigr)
    \twoheadrightarrow
    \bigl(\Omega^\bullet_c(I_1)\widehat\otimes\bar A\bigr)\oplus
    \bigl(\Omega^\bullet_c(I_2)\widehat\otimes\bar A\bigr)$$
  obtained by killing the diagonal multi-trace component (a tensor
  product of two distinct linear primitives is multi-trace and is
  $\pi_{\mathrm{prim}}$-zero on the indecomposable-quotient target).
\end{thm}

\begin{proof}
  \emph{Top-down route.}  A typical class on $\mathrm{Obs}^{\mathrm{cl}}_{\mathrm{open},N}(I_1)
  \otimes\mathrm{Obs}^{\mathrm{cl}}_{\mathrm{open},N}(I_2)$ is a tensor
  $\alpha\otimes\beta$ where each factor lies in some
  $F_{k_j}/F_{k_j-1}$.  The brane factorization product $\mu$ extends
  $\alpha\otimes\beta$ by zero from $I_j$ to $V$ and multiplies; the
  result lies in $F_{k_1+k_2}/F_{k_1+k_2-1}$.

  Apply $\pi_{\mathrm{prim}}$.  The projection kills everything in
  $F_{\geq 2}$ (Lemma~\ref{lem:cumulant-kernel}).  Therefore the
  output is nonzero only when $k_1+k_2\leq 1$.  The pairs $(k_1,k_2)$
  satisfying $k_1+k_2\leq 1$ are $(0,0)$, $(1,0)$, $(0,1)$.  The
  $(0,0)$ pair is killed by passing to $\bar A=\C[z_1,z_2]/\C\cdot1$
  (constants are excluded from the Hamiltonian quotient).  The
  $(1,0)$ pair sends a single-trace $B_w(a)$ on $I_1$ tensored with a
  scalar de Rham form on $I_2$ to the single-trace $B_w(a)$ on $V$,
  multiplied by the same scalar form on $I_2$, and similarly for
  $(0,1)$.  Modulo the scalar de Rham form, this is exactly the
  cosheaf extension-by-zero map $\mathfrak h_{I_j}\hookrightarrow
  \mathfrak h_V$ for $j=1,2$.

  \emph{Bottom-up route.}  $\pi_{\mathrm{prim}}\otimes\pi_{\mathrm{prim}}$
  applied to $\alpha\otimes\beta$ produces nonzero output only when
  each factor is in $F_1/F_0$.  When both factors are nonzero linear
  primitives, the augmentation map kills the tensor (this is the
  multi-trace-vanishing in the indecomposable target,
  Remark~\ref{rmk:cumulant-target-indecomposable}); when only one is
  nonzero (and the other is the unit), the augmentation is the
  identity on that factor.  The cosheaf extension by zero
  $\mu_{\bar A}$ then sends this to its image in
  $\Omega^\bullet_c(V)\widehat\otimes\bar A$.

  Both routes give the same answer in every stratum: the linear
  cosheaf extension-by-zero map of the brane $P_0$ factorization
  algebra at the indecomposable level.  Hence the diagram commutes
  strictly with $h_{\mathrm{fact}}\equiv0$.

  This is the chain-level statement of cosheaf compatibility.
\end{proof}

\begin{thm}[\(P_0\) bracket compatibility, with explicit homotopy]\label{thm:chain-level-primitive-projection-P0}
  For $F,G\in\mathrm{Obs}^{\mathrm{cl}}_{\mathrm{open},N}(I)$ in the
  stable connected limit,
  $$\pi_{\mathrm{prim}}\bigl(\{F,G\}_{\mathrm{open}}\bigr)
    -\bigl\{\pi_{\mathrm{prim}}F,\pi_{\mathrm{prim}}G\bigr\}_{\bar A}
    =(Q\circ h_{P_0}+h_{P_0}\circ Q)(F\otimes G),$$
  with the explicit homotopy $h_{P_0}\equiv 0$ on every stratum
  $F_{k_1}\otimes F_{k_2}$. The chain-level $P_0$-bracket compatibility
  square commutes strictly.
\end{thm}

\begin{proof}
  We compute on each stratum separately.

  \emph{Stratum $(k_1,k_2)=(1,1)$.}  Both inputs are single traces
  $B_f(a)$ and $B_g(b)$ for $f,g\in\bar A$ and
  $a,b\in\Omega^0_c(I)$.  By
  Lemma~\ref{lem:filtration-preservation}, the open
  $P_0$ bracket of two single traces is again a single trace
  (cycle merging in Procesi-Razmyslov), and by
  Theorem~\ref{thm:bulk-boundary}, in the connected Hamiltonian
  quotient,
  $\{B_f(a),B_g(b)\}_{\mathrm{open}}=B_{\{f,g\}_{\bar A}}(ab)$.
  Apply $\pi_{\mathrm{prim}}$: the result is the single-trace
  $ab\otimes\{f,g\}_{\bar A}$ in
  $\Omega^\bullet_c(I)\widehat\otimes\bar A$.

  Independently, $\pi_{\mathrm{prim}}B_f(a)=a\otimes f$ and
  $\pi_{\mathrm{prim}}B_g(b)=b\otimes g$.  The bracket on
  $\Omega^\bullet_c(I)\widehat\otimes\bar A$ is the
  $P_0$ bracket of Construction~\ref{constr:interval-fact-algebras}:
  $\{a\otimes f,b\otimes g\}_{\bar A}=ab\otimes\{f,g\}_{\bar A}$.
  This matches.  Strict compatibility on this stratum.

  \emph{Stratum $(k_1,k_2)=(1,k)$ with $k\geq 2$.}  $F=B_f(a)$ is a
  single trace; $G$ is a $k$-trace product
  $B_{w_1}(b_1)\cdots B_{w_k}(b_k)$.  By the Leibniz rule,
  $$\{F,G\}_{\mathrm{open}}=\sum_{i=1}^k(\pm)
    B_{w_1}(b_1)\cdots
    \{B_f(a),B_{w_i}(b_i)\}_{\mathrm{open}}\cdots
    B_{w_k}(b_k).$$
  Each term is a $k$-trace product (cycle merging joined to a
  remaining $(k-1)$-trace tail).  Hence the result lies in $F_k$,
  and by Lemma~\ref{lem:cumulant-kernel}, $\pi_{\mathrm{prim}}$ of
  this is zero whenever $k\geq 2$.

  Independently, $\pi_{\mathrm{prim}}G=0$ because $G\in F_{\geq 2}$,
  hence $\{\pi_{\mathrm{prim}}F,\pi_{\mathrm{prim}}G\}_{\bar A}=0$.

  Both sides are zero on this stratum.  Strict.

  \emph{Stratum $(k_1,k_2)$ with $k_1\geq 2$ or $k_2\geq 2$.}  By
  symmetry the same argument: both sides vanish.

  \emph{Stratum $(0,k)$ or $(k,0)$.}  One factor is in
  $F_0$, that is, a scalar de Rham form.  The bracket of a scalar
  de Rham form with anything is zero (the canonical equal-time
  bracket only acts on matrix fields).  So
  $\{F,G\}_{\mathrm{open}}=0$ and
  $\{\pi_{\mathrm{prim}}F,\pi_{\mathrm{prim}}G\}_{\bar A}=0$.  Strict.

  In every stratum the difference is zero; setting $h_{P_0}\equiv 0$
  realizes the homotopy formula trivially.
\end{proof}

\begin{rmk}[Why the homotopy is zero]\label{rmk:why-homotopy-zero}
  The strictness in
  Theorem~\ref{thm:chain-level-primitive-projection-P0} is not an
  accident of small word degree; it is a structural consequence of
  the trace-count filtration of
  Lemma~\ref{lem:filtration-preservation}.  The bracket of two single
  traces is a single trace by Procesi--Razmyslov cycle merging; the
  bracket of a single trace with a multi-trace is a multi-trace by
  the Leibniz rule; cumulant projection projects to the single-trace
  component.  Hence the failure of the bracket compatibility, if any,
  must lie in the multi-trace stratum, where both sides are zero.
  \end{rmk}

\begin{constr}[Primitive one-antifield chain complex]\label{constr:T5-one-psi-complex}
  Put \(x=\phi_1\), \(y=\phi_2\), and let
  \(\mathsf W_{k,l}\) be the set of ordinary words with \(k\) letters
  \(x\) and \(l\) letters \(y\).  Let \(\mathsf N_{a,b}\) be the set of
  cyclic words with \(a\) letters \(x\) and \(b\) letters \(y\).  The
  primitive one-\(\psi\) chain model in bidegree \((k,l)\) is
  \[
    C_2(k,l)\xrightarrow{\partial_2}C_1(k,l)
      \xrightarrow{\partial_1}C_0(k,l),
  \]
  where
  \[
    C_1(k,l)=\C\{\mathsf W_{k,l}\},\qquad
    C_0(k,l)=\C\{\mathsf N_{k+1,l+1}\}.
  \]
  The basis vector \(w\in\mathsf W_{k,l}\) represents
  \(\operatorname{Tr}(\psi w)\), with the graded cyclic rule
  \(\operatorname{Tr}(ab)=(-1)^{|a||b|}\operatorname{Tr}(ba)\),
  \(|x|=|y|=0\), \(|\psi|=1\).  The first boundary is
  \[
    \partial_1(w)=[xyw]_{\mathrm{cyc}}-[yxw]_{\mathrm{cyc}}.
  \]
  The two-\(\psi\) source is the alternating quotient generated by pairs
  of ordinary words \(u,v\), of total bidegree \((k-1,l-1)\), with
  \[
    [u,v]=-[v,u],\qquad [u,u]=0.
  \]
  In this quotient
  \[
    C_2(k,l)=
    \bigl(\C\{(u,v):\deg u+\deg v=(k-1,l-1)\}\bigr)_{\mathrm{alt}},
  \]
  and \(C_2(k,l)=0\) if \(k=0\) or \(l=0\).
  Thus the displayed complex is literally the cellular chain complex in
  dimensions \(2,1,0\): \(C_0\) is the necklace vertex set,
  \(C_1\) is the one-moving-chip edge set, and \(C_2\) is the oriented
  two-moving-chip square set.
\end{constr}

\begin{lem}[Two-antifield boundary sign]\label{lem:T5-two-psi-boundary-sign}
  In the chain model of
  Construction~\ref{constr:T5-one-psi-complex}, the boundary of the
  generator represented by \(\operatorname{Tr}(\psi u\psi v)\) is
  \[
    \partial_2(u,v)=vxyu-vyxu-uxyv+uyxv.
  \]
  This formula is well-defined on the alternating quotient and satisfies
  \(\partial_1\partial_2=0\).
\end{lem}

\begin{proof}
  The differential is the odd derivation determined by
  \(Q\psi=xy-yx\), \(Qx=Qy=0\).  Since the second \(\psi\) is preceded
  by one odd letter, replacing it carries the Koszul sign \(-1\):
  \[
  \begin{aligned}
    Q\,\operatorname{Tr}(\psi u\psi v)
      &=
      \operatorname{Tr}\bigl((xy-yx)u\psi v\bigr)
      -\operatorname{Tr}\bigl(\psi u(xy-yx)v\bigr).
  \end{aligned}
  \]
  In the first term the even word \((xy-yx)u\) is moved cyclically past
  the final displayed \(\psi v\), producing no additional sign.  Removing
  the common front \(\psi\) gives
  \(vxyu-vyxu-uxyv+uyxv\).  Interchanging \(u\) and \(v\) reverses this
  expression:
  \[
    uxyv-uyxv-vxyu+vyxu=-(vxyu-vyxu-uxyv+uyxv),
  \]
  so the formula descends to the alternating quotient \(C_2(k,l)\).

  Applying \(\partial_1\) to the four summands gives the signed sum of
  cyclic words obtained by inserting two adjacent commutators.  Each
  final cyclic word occurs twice with opposite signs, according to the
  order in which the two displayed \(\psi\)'s are replaced.  Equivalently
  \(Q^2=0\) for the odd derivation \(Q\).  Hence
  \(\partial_1\partial_2=0\).
\end{proof}

\begin{thm}[Primitive one-antifield homology]\label{thm:T5-one-psi-homology}
  For every \(k,l\ge0\),
  \[
    H_1\bigl(C_2(k,l)\to C_1(k,l)\to C_0(k,l)\bigr)
      \cong \C\cdot[\Psi_{k,l}].
  \]
  The Hamiltonian comparison uses the nonconstant sum \(k+l>0\); the
  scalar class \(\Psi_{0,0}=\operatorname{Tr}(\psi)\) is not part of
  \(\bar A_{\mathrm{desc}}\).
\end{thm}

\begin{proof}
  If \(k=0\) or \(l=0\), then \(C_1(k,l)\) is spanned by the single word
  \(x^k\) or \(y^l\), \(\partial_1=0\), and \(C_2(k,l)=0\).  The class is
  represented by \(\Psi_{k,l}\).

  Assume \(k,l>0\).  Form the graph \(\Gamma_{k,l}\) with vertices
  \(\mathsf N_{k+1,l+1}\) and oriented edges \(w\in\mathsf W_{k,l}\),
  \[
    w:[yxw]_{\mathrm{cyc}}\longrightarrow[xyw]_{\mathrm{cyc}}.
  \]
  Its cellular boundary is \(\partial_1\).  Attach one square for each
  alternating pair \((u,v)\) with boundary
  \(vxyu-vyxu-uxyv+uyxv\).  By
  Lemma~\ref{lem:T5-two-psi-boundary-sign}, this square boundary is
  exactly \(\partial_2\); denote the resulting two-complex by
  \(X_{k,l}\).

  Put \(R=k+1\) and \(S=l+1\).  \(X_{k,l}\) is the two-skeleton of the
  cyclic chip-moving state complex \(Y_{R,S}\).  Before cyclic
  relabelling, a \(p\)-cell is a configuration of \(R-p\) stationary
  indistinguishable chips on the vertices of an \(S\)-edge oriented
  circle and \(p\) moving chips in distinct edge interiors.  The
  orientation is the increasing order of the moving edge labels.  Reading
  around the labelled circle writes \(x\) for stationary chips, \(y\) for
  an empty edge, and \(\psi\) for a moving edge.  Hence vertices are
  \(\mathsf N_{R,S}\), one-cells are the words \(w\in\mathsf W_{k,l}\),
  and two-cells are exactly the alternating pairs \((u,v)\) above.

  The Abel map sends a labelled configuration to the sum of its chip
  positions in \(\R/S\Z\) and is affine on every cell.  After lifting the
  target to \(\R\) and fixing a cyclic order of the lifted chips, a fibre
  is the convex polytope cut out by
  \[
    t_1\leq\cdots\leq t_R\leq t_1+S,\qquad \sum_i t_i=\mathrm{constant},
  \]
  together with the chip-cell inequalities.  The chip-moving cells
  subdivide this polytope.  Thus the labelled Abel map is a homology
  equivalence with \(S^1\).

  Cyclic relabelling by \(r\in C_S\) adds \(r\) to every chip position, so
  the Abel coordinate is translated by \(Rr\).  The action on the target
  circle is orientation-preserving and therefore trivial on \(H_0\) and
  \(H_1\).  Since \(\C[C_S]\) is semisimple, homology commutes with
  coinvariants.  If a stabilizer \(G\subset C_S\) fixes an Abel value, it
  acts affinely on each lifted convex fibre; averaging gives a
  \(G\)-fixed point and straight-line contraction to it descends to a
  contraction of the quotient fibre.  Therefore the cyclic quotient still
  has the homology of \(S^1\).  Finally, attaching cells of dimension
  at least three cannot change \(H_1\), because \(H_1\) only sees
  \(C_2\to C_1\to C_0\).  Hence \(H_1(X_{k,l};\C)\cong\C\).

  The symmetrised trace
  \[
    \Psi_{k,l}=\sum_{w\in\mathsf W_{k,l}}\operatorname{Tr}(\psi w)
  \]
  is a cycle: the coefficient of any zero-\(\psi\) necklace in
  \(Q\Psi_{k,l}\) counts cyclic \(xy\)-transitions minus cyclic
  \(yx\)-transitions, and these two numbers are equal on a circle.  The
  functional \(\epsilon(\operatorname{Tr}(\psi w))=1\) kills every
  two-\(\psi\) boundary, because the formula in
  Lemma~\ref{lem:T5-two-psi-boundary-sign} has two positive and two
  negative summands, while
  \[
    \epsilon(\Psi_{k,l})=\binom{k+l}{k}\neq0.
  \]
  Each one-cell occurring in \(\Psi_{k,l}\) is oriented in the positive
  Abel direction.  Hence its Abel image is
  \(\binom{k+l}{k}\) times the fundamental circle class.  The unique
  \(H_1\)-line is generated by \([\Psi_{k,l}]\).
\end{proof}

\begin{defn}[Full primitive \(\psi\)-complex]\label{def:T5-full-psi-complex}
  Let
  \[
    A_\psi=\C\langle x,y,\psi\rangle,\qquad
    |x|=|y|=0,\quad |\psi|=1,
  \]
  with \(d\psi=xy-yx\) and \(dx=dy=0\).  In the graded cyclic quotient
  \(\operatorname{Cyc}(A_\psi)=A_\psi/[A_\psi,A_\psi]_{\mathrm{gr}}\),
  put
  \[
    \operatorname{wt}(x)=(1,0),\qquad
    \operatorname{wt}(y)=(0,1),\qquad
    \operatorname{wt}(\psi)=(1,1).
  \]
  For fixed \((R,S)\), define
  \[
    K^p_{R,S}
      =
      \C\{[w]_{\mathrm{cyc}}:
        \#\psi(w)=p,\ \#x(w)=R-p,\ \#y(w)=S-p\},
  \]
  with \(0\le p\le\min(R,S)\).  The boundary is
  \[
  \begin{aligned}
    \partial[w]
      =
      \sum_{i:a_i=\psi}(-1)^{\#\{\psi\text{ before }i\}}
      [&a_1\cdots a_{i-1}xy\,a_{i+1}\cdots a_n  \\
       &-a_1\cdots a_{i-1}yx\,a_{i+1}\cdots a_n]_{\mathrm{cyc}} .
  \end{aligned}
  \]
  The sign is the Koszul sign for replacing the \(i\)-th odd letter by
  the even commutator \(xy-yx\); with this sign \(\partial^2=0\).
\end{defn}

\begin{lem}[All-\(\psi\) face sign]\label{lem:T5-full-p-face-sign}
  Let \(u_1,\ldots,u_p\) be ordinary words in \(x,y\).  In
  \(K^{p-1}_{R,S}\),
  \[
  \begin{aligned}
    \partial[\psi u_1\psi u_2\cdots\psi u_p]
      =
      \sum_{j=1}^p(-1)^{j-1}
      [&\psi u_1\cdots\psi u_{j-1}xy\,u_j
        \psi u_{j+1}\cdots\psi u_p\\
       &-\psi u_1\cdots\psi u_{j-1}yx\,u_j
        \psi u_{j+1}\cdots\psi u_p]_{\mathrm{cyc}} .
  \end{aligned}
  \]
  Under the chip-moving cellular model below, this is exactly the
  oriented cubical boundary of the corresponding \(p\)-cell.
\end{lem}

\begin{proof}
  The \(j\)-th displayed \(\psi\) has precisely \(j-1\) odd letters before
  it.  Replacing it by the even commutator \(xy-yx\) therefore contributes
  the Koszul sign \((-1)^{j-1}\) and no further sign.  The cellular
  \(p\)-cell is oriented by the increasing order of its \(p\) moving
  chip-edges.  Its \(j\)-th coordinate has terminal face \(xy\) and
  initial face \(yx\), with cubical sign \((-1)^{j-1}\).  These two
  formulae agree term by term.  If a cyclic rotation moves a block
  containing \(a\) odd letters past a block containing \(b\) odd letters,
  the cellular orientation changes by \((-1)^{ab}\), the graded cyclic
  sign.  Hence the formula descends to the cyclic quotient.
\end{proof}

\begin{constr}[Chip-moving cellular model]\label{constr:T5-chip-moving-cellular-model}
  Assume \(R,S>0\).  Let \(\widetilde Y_{R,S}\) be the labelled state
  complex on an oriented \(S\)-edge circle.  A \(p\)-cell consists of
  \(R-p\) stationary indistinguishable chips on vertices and \(p\)
  moving chips in the interiors of \(p\) distinct edges.  Reading around
  the labelled circle writes \(x\) for each stationary chip at a vertex
  and writes \(y\), respectively \(\psi\), for an empty, respectively
  moving, following edge.  Thus \(p\)-cells are labelled words with
  \(R-p\) \(x\)'s, \(S-p\) \(y\)'s, and \(p\) \(\psi\)'s.

  The cyclic group \(C_S\) rotates the labelled edge set.  The quotient
  chain complex is taken with the orientation character already encoded
  by the graded cyclic rule
  \([uv]=(-1)^{|u||v|}[vu]\).
\end{constr}

\begin{lem}[Cellular chain identification]\label{lem:T5-cellular-chain-identification}
  For \(R,S>0\), the cellular chains of
  Construction~\ref{constr:T5-chip-moving-cellular-model}, after cyclic
  relabelling, identify with the full primitive cyclic Koszul complex:
  \[
    C_p(\widetilde Y_{R,S};\C)_{C_S}\cong K^p_{R,S}.
  \]
  Under this identification the cellular boundary is the boundary of
  Definition~\ref{def:T5-full-psi-complex}.
\end{lem}

\begin{proof}
  A labelled \(p\)-cell has a unique expression
  \(x^{a_1}m_1\cdots x^{a_S}m_S\), where
  \(m_i\in\{y,\psi\}\), exactly \(p\) markers are \(\psi\), and
  \(\sum a_i=R-p\).  Every graded cyclic word admits such a representative
  after moving the cut to an edge marker; choosing the labelled starting
  edge makes the representative unique.  Forgetting the labelled starting
  edge is exactly passage to graded cyclic words, with no relation other
  than cyclic relabelling.  The sign check is
  Lemma~\ref{lem:T5-full-p-face-sign}; cyclic relabelling by a rotation
  that moves \(a\) moving-edge coordinates past \(b\) moving-edge
  coordinates changes orientation by \((-1)^{ab}\), which is the same
  sign as moving the corresponding odd \(\psi\)-blocks in the graded
  cyclic quotient.
\end{proof}

\begin{lem}[Abel map and cyclic stabilizers]\label{lem:T5-abel-cyclic-stabilizers}
  The quotient state complex
  \(Y_{R,S}=\widetilde Y_{R,S}/C_S\), with the orientation local system
  of Lemma~\ref{lem:T5-cellular-chain-identification}, has the homology
  of \(S^1\).  The surviving \(H_1\)-line is represented by the
  Abel-positive cycle \(\Psi_{R-1,S-1}\).
\end{lem}

\begin{proof}
  Work first before cyclic relabelling.  In the universal cover of the
  \(S\)-edge circle, a lifted configuration is described by weakly
  ordered chip positions
  \[
    t_1\leq\cdots\leq t_R\leq t_1+S.
  \]
  The Abel map sends the configuration to
  \(\sum_i t_i\in\R/S\Z\) and is affine on every cell.  After lifting the
  target circle to \(\R\), each fibre is the convex polytope cut out by
  the displayed inequalities and by the fixed-sum equation.  The
  chip-moving cells subdivide this polytope.  Hence every lifted fibre is
  a contractible finite polyhedral complex, and the labelled Abel map is a
  proper PL map with acyclic fibres.  Vietoris--Begle gives a cellular
  homology equivalence
  \[
    \widetilde Y_{R,S}\simeq_{\mathrm H} S^1 .
  \]

  A cyclic relabelling by \(r\in C_S\) adds \(r\) to each chip position,
  so the Abel coordinate is translated by \(Rr\) on \(\R/S\Z\).  Thus
  relabelling acts on the target by orientation-preserving rotations and
  acts trivially on \(H_0\) and \(H_1\).  Since \(\C[C_S]\) is
  semisimple, taking coinvariants commutes with homology:
  \[
    H_\bullet(C_\bullet(\widetilde Y_{R,S})_{C_S})
      \cong H_\bullet(\widetilde Y_{R,S})_{C_S}
      \cong H_\bullet(S^1).
  \]

  The fibre statement survives stabilizers.  If a subgroup \(G\subset C_S\)
  fixes a target Abel value, choose the lifted fibre before decomposing it
  into chip cells.  It is a convex polytope preserved by the affine
  \(G\)-action.  Averaging any point of the fibre gives a \(G\)-fixed point,
  and the straight-line contraction to that fixed point is \(G\)-equivariant
  and respects the chip-cell subdivision.  Therefore the stabilizer quotient
  of the fibre is contractible.

  Finally, \(\Psi_{R-1,S-1}=\sum_w[\psi w]_{\mathrm{cyc}}\) is closed by
  the same transition count used in
  Theorem~\ref{thm:T5-one-psi-homology}.  Each of its one-cells is
  oriented in the positive Abel direction, so its image is a nonzero
  multiple of the fundamental class of the target circle.  Hence it
  spans the unique \(H_1\)-line.
\end{proof}

\begin{thm}[Full primitive higher-\(\psi\) Koszul homology]
  \label{thm:T5-full-primitive-psi-homology}
  For every \(R,S\ge0\), the full primitive cyclic Koszul complex of
  Definition~\ref{def:T5-full-psi-complex} has homology
  \[
    H_p(K_{R,S})\cong
    \begin{cases}
      \C\cdot[x^Ry^S],&p=0,\\
      \C\cdot[\Psi_{R-1,S-1}],&p=1,\ R,S\ge1,\\
      0,&p\ge2.
    \end{cases}
  \]
  Equivalently, the stable primitive single-trace Koszul homology is
  \[
    \C[x,y]\oplus\C[x,y]\theta,\qquad
    |\theta|=1,\qquad \operatorname{wt}(\theta)=(1,1),
  \]
  with \(x^ay^b\theta\) represented by \(\Psi_{a,b}\).
\end{thm}

\begin{proof}
  If \(R=0\) or \(S=0\), no positive \(\psi\)-degree term exists and
  \(K^0_{R,S}\) is the one-dimensional span of the necklace \(x^R\) or
  \(y^S\), including the empty word when \(R=S=0\).

  Assume \(R,S>0\).  Lemma~\ref{lem:T5-cellular-chain-identification}
  identifies \(K^\bullet_{R,S}\) with the cyclic-coinvariant cellular
  chains of \(\widetilde Y_{R,S}\).  Lemma~\ref{lem:T5-abel-cyclic-stabilizers}
  computes their homology as \(H_\bullet(S^1;\C)\).  Thus only degrees
  \(0\) and \(1\) survive.  The \(H_0\)-line is the commutative necklace
  class \([x^Ry^S]\).  The \(H_1\)-line is the Abel-positive fundamental
  class, represented by \(\Psi_{R-1,S-1}\).
  The vanishing for \(p\ge2\) is the vanishing of \(H_{\ge2}(S^1;\C)\);
  the \(H_1\) calculation is already visible in the three-term skeleton,
  since higher cells do not alter
  \(\ker(C_1\to C_0)/\operatorname{im}(C_2\to C_1)\).
\end{proof}

\begin{thm}[Hamiltonian sector extracted from full psi homology]
  \label{thm:T5-hamiltonian-sector-extraction}
  The Hamiltonian comparison uses exactly the nonconstant summands
  \[
    \bar A\oplus \bar A_{\mathrm{desc}}[1]
      =
    \bigoplus_{R+S>0}\C\cdot[x^Ry^S]
      \oplus
    \bigoplus_{a+b>0}\C\cdot[\Psi_{a,b}].
  \]
  The sectors outside that comparison are precisely:
  the scalar degree-zero class \([1]\), the scalar odd class
  \(\Psi_{0,0}=\operatorname{Tr}(\psi)\), and the absent higher
  primitive \(\psi\)-sectors \(H_p(K_{R,S})=0\) for \(p\ge2\).
  Identifying the remaining \(\Psi_{a,b}\) line with principal parts
  \(\rho_{a,b}\) as an \(\mathfrak h\)-module, or after Moyal
  quantization, is not part of this homology theorem.
\end{thm}

\begin{proof}
  Constants are removed from the closed Hamiltonian algebra
  \(\bar A=\C[z_1,z_2]/\C\cdot1\), and the open empty trace is the scalar
  \(\operatorname{Tr}(1)=N\).  This removes the \(H_0(K_{0,0})\) line.
  The class \(\Psi_{0,0}\) is a genuine odd primitive class, but it is the
  descendant of the removed constant Hamiltonian and therefore has no
  closed Hamiltonian partner in \(\bar A_{\mathrm{desc}}[1]\).  By
  Theorem~\ref{thm:T5-full-primitive-psi-homology}, all primitive
  \(p\ge2\) homology groups vanish.  The only surviving positive
  Hamiltonian labels are therefore the displayed nonconstant \(H_0\) and
  \(H_1\) lines.  The final sentence is exactly the module-structure
  separation proved later by
  Proposition~\ref{prop:quantum-descendant-label-obstruction} and
  Theorem~\ref{thm:one-psi-label-hbar-obstruction}.
\end{proof}

\begin{thm}[Universal primitive descendant envelope]
  \label{thm:T5-universal-primitive-descendant-envelope}
  The stable primitive source for the reduced cotangent comparison is forced
  to be
  \[
    \bar A\oplus\bar A_{\mathrm{desc}}[1]
    =
    \bigoplus_{R+S>0}\C\cdot[x^Ry^S]
      \oplus
    \bigoplus_{a+b>0}\C\cdot[\Psi_{a,b}].
  \]
  More precisely, every stable primitive chain-level construction generated
  by \(x,y,\psi\), with \(Q\psi=xy-yx\) and with only graded-cyclic trace
  relations, factors on homology through the two displayed summands.  Thus
  a universal Koszul-dual cotangent theory in this lane has no available
  hidden higher-\(\psi\) correction term: the entire descendant side is the
  one-\(\psi\) line \(\bar A_{\mathrm{desc}}[1]\).

  After the Matlis normalization
  \[
    \langle\rho_{a,b},z_1^m z_2^n\rangle
      =\delta_{a,m}\delta_{b,n},
  \]
  the only flat \(D_\hbar\)-bridge constructed from the polynomial
  descendant labels is the Fourier-twisted Rees bridge of
  Theorem~\ref{thm:T5-fourier-rees-positive-construction}.  If the original
  Hamiltonian action of \(z_1,z_2\) is required to be preserved, the exact
  obstruction is the local-nilpotence mismatch of
  Theorem~\ref{thm:no-same-action-rees-fourier-descendants}.
\end{thm}

\begin{proof}
  The first assertion is Theorem~\ref{thm:T5-full-primitive-psi-homology}
  together with the scalar removal in
  Theorem~\ref{thm:T5-hamiltonian-sector-extraction}.  The cellular proof
  identifies the universal primitive complex with the chip-moving complex
  and proves \(H_p(K_{R,S})=0\) for \(p\ge2\).  Therefore any stable
  primitive construction satisfying the same trace and differential
  hypotheses has no homology in which an additional universal cotangent
  correction could live.

  The Matlis pairing fixes the dual coefficient normalization up to the
  single scalar allowed by residue-pairing rigidity; the displayed formula
  chooses that scalar.  With this normalization, the positive flat Rees
  construction is exactly the Fourier-twisted cyclic \(D_\hbar\)-module
  \(\mathcal D_\hbar/(z_1,z_2)\) embedded into the principal-part Cech
  lattice with the explicit factor
  \((-1)^{a+b}\hbar^{a+b}a!b!\).  If one removes the Fourier twist and asks
  for the same Hamiltonian action, local nilpotence of the linear
  Hamiltonians is preserved by flat generic-fibre isomorphism but differs
  on the two sides.  This is precisely the obstruction theorem cited in
  the statement.
\end{proof}

\begin{constr}[Primitive one-\(\psi\) label shadow]\label{constr:chain-level-primitive-cotangent}
  Extend $\pi_{\mathrm{prim}}$ to the one-$\psi$ primitive shadow by
  recording the labels of the symmetrised traces $\Psi_{a,b}$ of
  Construction~\ref{constr:boundary-evaluation}.  Define
  $$\pi_{\mathrm{prim}}^{\mathrm{cot}}:
    \mathrm{Obs}^{\mathrm{cl,cot}}_{\mathrm{open},N}(I)
    \longrightarrow
    \Omega^\bullet_c(I)\widehat\otimes\bigl(\bar A\oplus\mc P[1]\bigr)$$
  by:
  \begin{itemize}
    \item on the zero-$\psi$ Hamiltonian sector,
      $\pi_{\mathrm{prim}}^{\mathrm{cot}}=\pi_{\mathrm{prim}}$
      composed with the inclusion $\bar A\hookrightarrow\bar A\oplus\mc P[1]$;
    \item on a primitive one-$\psi$ class
      $\Psi_{a,b}$ smeared by a distributional test current
      $B\in\mathcal D'_c(I)$,
	      $$\pi_{\mathrm{prim}}^{\mathrm{cot}}\bigl(B\otimes\Psi_{a,b}\bigr)
	        =B\otimes\rho_{a,b}[1]
	        =\Theta_{a,b}(B);$$
    \item on multi-trace products of one-$\psi$ classes (or mixed
      products of zero-$\psi$ and one-$\psi$ traces),
      $\pi_{\mathrm{prim}}^{\mathrm{cot}}\equiv 0$ (cumulant on the
      multi-trace stratum).
  \end{itemize}
	  This is a shadow map to the reduced principal-part label target, not
	  a polynomial realization of the coadjoint module. The
	  no-polynomial theorem \ref{thm:no-polynomial-realization-categorical}
	  says precisely that the open polynomial one-$\psi$ descendants and
	  the principal-part currents do not carry the same
	  $\mathfrak h$-module structure. The map above is therefore an
	  indecomposable shadow map after imposing the reduced current target
	  $\mc P[1]$; it is not an unreduced boundary observable realizing
	  $\mc P$ inside the polynomial BV/Koszul algebra.
\end{constr}

\begin{thm}[Primitive indecomposable \(P_0\)-shadow and one-\(\psi\) label shadow]\label{thm:chain-level-primitive-projection}
  In the stable connected limit, the primitive indecomposable \(P_0\)-shadow
  $\pi_{\mathrm{prim}}$ of
  Construction~\ref{constr:chain-level-primitive-projection} (and its
  one-$\psi$ shadow $\pi_{\mathrm{prim}}^{\mathrm{cot}}$ of
  Construction~\ref{constr:chain-level-primitive-cotangent}) is a
  surjection to the indecomposable $P_0$ shadow of the boundary
  factorization algebra
  $$\pi_{\mathrm{prim}}^{\mathrm{cot}}:
    \widetilde{\mathrm{Obs}}^{\mathrm{cl,pp}}_{\partial,N}
    \twoheadrightarrow
    A_\partial^{\mathrm{pp}}$$
  in the sense of Construction~\ref{constr:interval-fact-algebras},
  with the following compatibilities:
  \begin{enumerate}
    \item \emph{(BV chain map.)}
      $\pi_{\mathrm{prim}}^{\mathrm{cot}}\circ Q=Q\circ\pi_{\mathrm{prim}}^{\mathrm{cot}}$
      strictly at the chain level
      (Theorem~\ref{thm:chain-level-primitive-projection-Q} and
      Theorem~\ref{thm:T5-one-psi-homology} extending it to the
      cotangent sector via the closed cocycle property of $\Psi_{a,b}$);
    \item \emph{(Factorization product.)}
      The factorization-product compatibility square commutes strictly
      modulo a zero homotopy
      (Theorem~\ref{thm:chain-level-primitive-projection-factorization}),
	      with the same trace-count argument on the one-$\psi$ label
	      projection;
		    \item \emph{($P_0$ bracket.)} The $P_0$-bracket compatibility
		      square commutes strictly modulo a zero homotopy on the
		      Hamiltonian zero-$\psi$ indecomposable quotient
		      (Theorem~\ref{thm:chain-level-primitive-projection-P0}). On the
		      one-$\psi$ shadow the map records labels only; it does
	      not intertwine the open polynomial descendant raising action
	      with the principal-part coadjoint lowering action. The reduced
	      target bracket
	      $\{B_f(a),\Theta_\rho(B)\}=\Theta_{f\cdot\rho}(aB)$ is supplied
	      separately by
	      Theorem~\ref{thm:reduced-principal-part-boundary-current};
    \item \emph{(Cosheaf.)} $\pi_{\mathrm{prim}}^{\mathrm{cot}}$ is
      compatible with extension by zero in $\Omega^\bullet_c$,
      because both source and target are tensored over the de Rham
      cosheaf $\Omega^\bullet_c$ on $\R$.
  \end{enumerate}
  The target is the nonunital indecomposable quotient: multiplication
  of two positive primitive generators is zero after projection,
  while the Lie/$P_0$ bracket remains the Hamiltonian bracket.  Thus
  the statement is not an algebra map to $\bar A$ with its ordinary
  commutative multiplication.

  Equivalently, $\pi_{\mathrm{prim}}^{\mathrm{cot}}$ is the chain-level
  realization of the unreduced-to-reduced primitive shadow required by
  Problem~\ref{prob:formal-factorization-center}, item~(3). It is not a
  proof of the unreduced cotangent factorization-center lift.
\end{thm}

\begin{proof}
  Items (1), (2), (3) on the Hamiltonian zero-$\psi$ sector are
  Theorems~\ref{thm:chain-level-primitive-projection-Q},
  \ref{thm:chain-level-primitive-projection-factorization},
  \ref{thm:chain-level-primitive-projection-P0} above.

  One-$\psi$ shadow.  By Theorem~\ref{thm:T5-one-psi-homology},
  $Q\Psi_{a,b}=0$ at the chain level.  Hence
  $Q\,\pi_{\mathrm{prim}}^{\mathrm{cot}}(B\otimes\Psi_{a,b})
  =\partial B\otimes\rho_{a,b}[1]=\pi_{\mathrm{prim}}^{\mathrm{cot}}(\partial B\otimes\Psi_{a,b}
   +B\otimes Q\Psi_{a,b})=\pi_{\mathrm{prim}}^{\mathrm{cot}}\circ Q(B\otimes\Psi_{a,b})$
  (the only contribution comes from the de Rham differential $\partial$
  on the distributional test current $B$, since $Q$ kills
	  $\Psi_{a,b}$).  The factorization-product compatibility on
	  one-$\psi$ generators uses the same trace-count
	  filtration: a one-$\psi$ trace tensored with another trace is a
	  multi-trace, killed by cumulant projection.  The
	  $P_0$-bracket compatibility asserted here is the zero-$\psi$
	  Hamiltonian statement already proved above. The one-$\psi$ shadow
	  deliberately does not assert compatibility between the open
	  descendant raising action and the principal-part coadjoint lowering
	  action; their difference is
	  Corollary~\ref{cor:descendant-coadjoint-difference}. The
	  bracket of two cotangent generators is zero on both sides because
	  $\mc P[1]$ is an abelian odd ideal in
  $\mathfrak g_I=\mathfrak h_I\ltimes\mathcal P_I[1]$.

  Cosheaf compatibility is automatic: extension by zero on the
  $\Omega^\bullet_c$ factor commutes with the projection on the
  trace-algebra factor because the projection is of the form
  $\mathrm{id}\otimes p$ with $p$ a fixed linear map on the trace
  algebra.

  The chain-level surjectivity is the surjectivity of the cumulant
  projection onto generators in the stable symmetric algebra (each
  generator is hit by itself), tensored with the identity on the
  compactly supported de Rham/current factor on the same interval.

  This establishes all four items.
\end{proof}

\begin{prop}[Quantum obstruction to the primitive one-\(\psi\) label map]
  \label{prop:quantum-descendant-label-obstruction}
  Let \(\mathfrak h_\hbar=\bar A_\hbar\) carry the Weyl/Moyal bracket
  \[
    \{f,g\}_\hbar
    =
    \sum_{j\geq0}
    \frac{\hbar^{2j}}{(2j+1)!\,2^{2j}}P^{2j+1}(f,g),
    \qquad
    P=\partial_{z_1}\otimes\partial_{z_2}
      -\partial_{z_2}\otimes\partial_{z_1}.
  \]
  The chain-level one-\(\psi\) primitive label map
  \(\Psi_{a,b}\mapsto\rho_{a,b}\) is not compatible with the
  \(\mathfrak h_\hbar\)-coadjoint action.  The failure already occurs in
  the indecomposable target; it is not a decomposable multi-trace
  correction.
\end{prop}

\begin{proof}
  Compatibility would require the polynomial PBW descendant action
  of Proposition~\ref{prop:open-descendant-action}, after the label
  replacement \(\widetilde\Psi_{a,b}\mapsto\rho_{a,b}\), to agree with
  the Moyal coadjoint action
  \[
    (\operatorname{ad}^{\vee}_{f,\hbar}\delta)(g)
    =-\delta(\{f,g\}_\hbar).
  \]
  Test this on the single Hamiltonian \(f=z_1^4\) and the single
  descendant label \((a,b)=(1,1)\).

  On the polynomial one-\(\psi\) side,
  Proposition~\ref{prop:open-descendant-action} gives
  \[
    z_1^4:\widetilde\Psi_{1,1}
    \longmapsto 4\,\widetilde\Psi_{4,0}.
  \]
  On the Moyal coadjoint side the classical Poisson term has negative
  first target index and vanishes, but the \(P^3\) term does not:
  \[
    P^3(z_1^4,z_2^4)=576\,z_1z_2,\qquad
    \frac{1}{3!\,2^2}P^3(z_1^4,z_2^4)=24\,z_1z_2.
  \]
  Hence
  \[
    \operatorname{ad}^{\vee}_{z_1^4,\hbar}\delta_{1,1}
    =-24\,\hbar^2\delta_{0,4}+\cdots,
  \]
  equivalently
  \(z_1^4\cdot_\hbar\rho_{1,1}=-24\,\hbar^2\rho_{0,4}+\cdots\) in the
  principal-part realization.  The two answers have different
  \(\hbar\)-order and different primitive label.  Since both sides have
  already been projected to the linear indecomposable summand, adding
  decomposable multi-trace terms cannot repair this square after
  applying \(\pi_{\mathrm{prim}}^{\mathrm{cot}}\).

	  The finite exact-arithmetic attack verifies the same
	  obstruction through exponents \(\leq5\) and Moyal order \(\leq5\):
  \(1200\) of \(1225\) nonzero test cases disagree under the naive label
  map, and \(958\) cases have genuine higher Moyal coadjoint terms.  The
  same finite computation also checks the corrected principal-part target below in
  \(3375\) finite representation identities with exponents \(\leq3\),
  Moyal order \(\leq5\), and \(\hbar\)-power \(\leq4\), and checks the
	  linear-generator leading-term obstruction to every \(\hbar\)-adic
	  deformation of the classical label shadow in \(70\) tests with
  labels \(\leq5\).
  The displayed \(z_1^4\) calculation is already a proof that the
  classical primitive one-\(\psi\) projection does not itself quantize
  to the descendant coadjoint sector.
\end{proof}

\begin{constr}[Reduced Moyal principal-part descendant target]
  \label{constr:quantum-principal-part-descendant-target}
  The corrected quantum descendant target is obtained by changing the
  target, not by correcting the polynomial one-\(\psi\) representatives.
  For an interval \(I\subset\R\) put
  \[
    \mathfrak h_{\hbar,I}
      =\Omega^\bullet_c(I)\widehat\otimes\bar A_\hbar,\qquad
    \mathcal P_{\hbar,I}
      =\mathcal D'_c(I)\widehat\otimes\mc P[[\hbar]].
  \]
  Define the reduced quantum principal-part descendant complex
  \[
    A^{\mathrm{pp}}_{\partial,\hbar}(I)
    =
    \widehat{\mathbf S}_{\C[[\hbar]]}(\mathfrak h_{\hbar,I})
    \widehat\otimes
    \widehat{\mathbf S}_{\C[[\hbar]]}(\mathcal P_{\hbar,I}[1]).
  \]
  The differential is the de Rham/current differential on the interval
  factors and zero on the coefficient modules \(\bar A_\hbar\) and
  \(\mc P[[\hbar]]\).  The bracket is the Moyal semidirect bracket
  \[
    \{B_f(a),B_g(b)\}_\hbar
      =B_{\{f,g\}_\hbar}(ab),\qquad
    \{B_f(a),\Theta_\rho(B)\}_\hbar
      =\Theta_{\operatorname{ad}^{\vee}_{f,\hbar}\rho}(aB),
  \]
  with
  \[
    \{\Theta_\rho(B),\Theta_\sigma(C)\}_\hbar=0.
  \]
  Here \(a,b\in\Omega^0_c(I)\), \(B,C\in\mathcal D'_c(I)\), and
  \(aB\) denotes multiplication of a compactly supported distribution by
  a smooth test function.

  In monomial coordinates the full coadjoint action is as follows.  Put
  \[
    C_r(p,q;m,n)=
      \sum_{\alpha=0}^{r}(-1)^\alpha\binom r\alpha
      (p)_{r-\alpha}(q)_\alpha(m)_\alpha(n)_{r-\alpha}.
  \]
  Then
  \[
    z_1^pz_2^q\cdot_\hbar\rho_{a,b}
    =
    -\sum_{\substack{r\geq1\\ r\ \mathrm{odd}}}
      \frac{\hbar^{r-1}}{2^{r-1}r!}\,
      C_r(p,q;a-p+r,b-q+r)\,
      \rho_{a-p+r,b-q+r},
  \]
  with the convention that negative indices and the scalar target
  \((0,0)\) are zero.  This is a finite sum for each monomial
  \(z_1^pz_2^q\), because the falling factorials vanish for
  \(r>p+q\).  The \(r=1\) term is the classical Matlis coadjoint action
  of Theorem~\ref{thm:principal-part-coadjoint-uniqueness}; the
  \(r\geq3\) terms are the genuine Moyal corrections.

  This construction supplies a reduced distributional principal-part
  target for the quantum descendant sector.  It does not make the
  polynomial PBW classes \(\Psi_{a,b}\) into that target.
\end{constr}

\begin{thm}[\(\hbar\)-adic obstruction for the one-\(\psi\) label shadow]
  \label{thm:one-psi-label-hbar-obstruction}
  Let \(M_{\Psi,\hbar}\) be a \(\C[[\hbar]]\)-flat complex generated
  over \(\C[[\hbar]]\) by lifts of the primitive polynomial
  one-\(\psi\) classes \(\widetilde\Psi_{a,b}\), \(a+b>0\).  Assume:
  \begin{enumerate}
    \item the source action of the linear Hamiltonians reduces modulo
      \(\hbar\) to the PBW descendant action
      \[
        z_1\cdot_0\widetilde\Psi_{a,b}
          =b\,\widetilde\Psi_{a,b-1},\qquad
        z_2\cdot_0\widetilde\Psi_{a,b}
          =-a\,\widetilde\Psi_{a-1,b};
      \]
    \item the target coefficient module is \(\mc P[[\hbar]][1]\),
      equivalently any fixed interval-current summand of
      \(\mathcal P_{\hbar,I}[1]\) from
      Construction~\ref{constr:quantum-principal-part-descendant-target},
      with the full Moyal coadjoint action;
    \item a candidate comparison
	      \(T_\hbar:M_{\Psi,\hbar}\to\mc P[[\hbar]][1]\) has the
	      classical one-\(\psi\) label shadow as leading term:
      \[
        T_\hbar(\widetilde\Psi_{a,b})
          =\rho_{a,b}[1]+O(\hbar).
      \]
  \end{enumerate}
  Then \(T_\hbar\) cannot intertwine the
  \(\mathfrak h_\hbar\)-action.  The obstruction already appears after
  reduction modulo \(\hbar\) and for the linear Hamiltonian \(z_1\).
  Consequently no strict quantum descendant lift of the classical
  primitive label shadow exists under these assumptions.  The only
  surviving quantum target in this lane is the reduced
  principal-part complex itself, not the polynomial one-\(\psi\) source.
\end{thm}

\begin{proof}
  Reduce a hypothetical intertwiner modulo \(\hbar\).  For the class
  \(\widetilde\Psi_{1,0}\), the PBW source action gives
  \[
    z_1\cdot_0\widetilde\Psi_{1,0}=0.
  \]
  Hence the source-then-target route gives
  \[
    T_0(z_1\cdot_0\widetilde\Psi_{1,0})=0.
  \]
  The target-then-action route gives, by the \(r=1\) term of
  Construction~\ref{constr:quantum-principal-part-descendant-target},
  \[
    z_1\cdot_0T_0(\widetilde\Psi_{1,0})
      =z_1\cdot\rho_{1,0}[1]
      =-\rho_{1,1}[1].
  \]
  The principal-part basis vector \(\rho_{1,1}\) is nonzero in
  \(\mc P\).  This contradiction is independent of all higher
  \(\hbar\)-corrections to \(T_\hbar\), because those corrections vanish
  after taking the leading term.  It is also independent of the higher
  Moyal terms: the linear Hamiltonian has no \(r\geq3\) Moyal
  correction.

  The same computation with \(z_2\) gives
  \(z_2\cdot_0\widetilde\Psi_{0,1}=0\) on the PBW side but
  \(z_2\cdot\rho_{0,1}=\rho_{1,1}\) on the principal-part side.  Thus
  the mismatch is not a special feature of one coordinate.
  Since the coefficient differential on \(\mc P[[\hbar]]\) is zero,
  this leading coefficient is not a boundary in the reduced target.
  A higher homotopy can remove it only after changing the target
  complex or the classical source action, which is outside the stated
  assumptions.
\end{proof}

\begin{thm}[Universal Fourier-twisted Rees D-hbar construction]
  \label{thm:T5-fourier-rees-positive-construction}
  Let
  \[
    D_\hbar=
    \C[[\hbar]]\langle z_1,z_2,\partial_1,\partial_2\rangle/
    ([\partial_i,z_j]=\hbar\delta_{ij},\ [z_i,z_j]=[\partial_i,\partial_j]=0)
  \]
  be the Rees Weyl algebra.  Let
  \[
    \mathcal O_\hbar=D_\hbar/(\partial_1,\partial_2),\qquad
    \Delta^{\mathsf F}_\hbar=D_\hbar/(z_1,z_2),
  \]
  and let \(e_0\) be the image of \(1\) in
  \(\Delta^{\mathsf F}_\hbar\).  The Fourier automorphism
  \[
    \mathsf F(z_i)=\partial_i,\qquad
    \mathsf F(\partial_i)=-z_i
  \]
  identifies the Fourier twist of \(\mathcal O_\hbar\) with
  \(\Delta^{\mathsf F}_\hbar\).  In the Cech local-cohomology lattice
  generated by
  \(\rho_{a,b}=z_1^{-a-1}z_2^{-b-1}dz_1dz_2\), the map
  \(e_0\mapsto\rho_{0,0}\) sends
  \[
    \partial_1^a\partial_2^b e_0
      \longmapsto
      (-1)^{a+b}\hbar^{a+b}a!b!\,\rho_{a,b}.
  \]
  This is a flat \(D_\hbar\)-module bridge in the Fourier-conjugated
  sense: \(P\in D_\hbar\) acts on the polynomial source as \(P\), and on
  the delta target as \(\mathsf F(P)\).  After tensoring with
  \(\C((\hbar))\), it gives a filtered Fourier--Rees label bijection after
  this polarization change between the nonconstant polynomial label span and
  \(\mc P((\hbar))\).  It is not an \(\mathfrak h\)-module identification
  for the original Hamiltonian action.  Before \(\hbar\) is inverted,
  the image is
  the explicit Rees sublattice
  \[
    \bigoplus_{a,b\ge0}
      \C[[\hbar]]\,(-1)^{a+b}\hbar^{a+b}a!b!\,\rho_{a,b}.
  \]
\end{thm}

\begin{proof}
  The formulas for \(\mathsf F\) preserve the Rees Weyl relation:
  \[
    [\mathsf F(\partial_i),\mathsf F(z_j)]
      =[-z_i,\partial_j]=\hbar\delta_{ij}.
  \]
  The polynomial module \(\mathcal O_\hbar\) is generated by \(1\)
  killed by \(\partial_1,\partial_2\).  After applying \(\mathsf F\), the
  cyclic generator is killed by \(z_1,z_2\), which is exactly the
  defining cyclic vector of \(\Delta^{\mathsf F}_\hbar\).

  In the Cech model \(\partial_i\) acts as
  \(\hbar\partial_{z_i}\).  Hence
  \[
  \begin{aligned}
    \partial_1^a\partial_2^b\rho_{0,0}
      &=
      \hbar^{a+b}
      \partial_{z_1}^a\partial_{z_2}^b
      (z_1^{-1}z_2^{-1}dz_1dz_2)  \\
      &=(-1)^{a+b}\hbar^{a+b}a!b!\rho_{a,b}.
  \end{aligned}
  \]
  The factorials and signs are units over \(\C\); the nontrivial Rees
  datum is the power \(\hbar^{a+b}\).  Inverting \(\hbar\) turns this
  sublattice into the whole principal-part label span.
\end{proof}

\begin{thm}[Exact same-action Rees obstruction]
  \label{thm:no-same-action-rees-fourier-descendants}
  There is no \(\C[[\hbar]]\)-flat Rees module \(M_\hbar\) with all three
  properties:
  \begin{enumerate}
    \item its special fibre is the polynomial primitive descendant module
      generated by the \(\widetilde\Psi_{a,b}\), \(a+b>0\), with linear
      Hamiltonian action
      \[
        z_1\cdot\widetilde\Psi_{a,b}=b\,\widetilde\Psi_{a,b-1},
        \qquad
        z_2\cdot\widetilde\Psi_{a,b}=-a\,\widetilde\Psi_{a-1,b};
      \]
    \item on the generic fibre \(M_\hbar((\hbar))\), the same linear
      Hamiltonians \(z_1,z_2\) act locally nilpotently, as they do in the
      polynomial Rees Weyl module;
    \item the generic fibre is isomorphic, as a module for the same
      Hamiltonian Lie algebra, to
      \(\mc P((\hbar))\) with the Matlis coadjoint action
      \[
        z_1\cdot\rho_{a,b}=-(b+1)\rho_{a,b+1},\qquad
        z_2\cdot\rho_{a,b}=(a+1)\rho_{a+1,b}.
      \]
  \end{enumerate}
  Equivalently, matching the local-nilpotence profile of \(z_1,z_2\) on
  the generic fibre is a necessary condition for any flat same-action
  Rees bridge from polynomial one-\(\psi\) descendants to principal parts,
  and that necessary condition fails.  Thus the universal bridge exists
  only after the Fourier twist of polarization in
  Theorem~\ref{thm:T5-fourier-rees-positive-construction}; it cannot be a
  same-action Rees family.
\end{thm}

\begin{proof}
  Local nilpotence of a fixed operator is invariant under module
  isomorphism over \(\C((\hbar))\).  On the polynomial descendant side,
  the displayed formulas lower the \(b\)-index for \(z_1\) and the
  \(a\)-index for \(z_2\); hence every finite vector is killed by a high
  enough power of each linear Hamiltonian.  The same is true for the
  polynomial Rees Weyl module after inverting \(\hbar\), because
  multiplying the operators by the invertible scalar \(\hbar\) does not
  change local nilpotence.

  On \(\mc P((\hbar))\), no nonzero finite Laurent vector is locally
  nilpotent.  Indeed, for
  \(v=\sum c_{a,b}\rho_{a,b}\neq0\),
  \[
    z_1^N v
      =
      \sum c_{a,b}(-1)^N(b+1)(b+2)\cdots(b+N)\rho_{a,b+N}.
  \]
  The basis vectors \(\rho_{a,b+N}\) remain distinct for distinct
  \((a,b)\), and every rising factorial displayed above is a nonzero
  element of \(\C((\hbar))\).  Thus \(z_1^Nv\neq0\) for all \(N\).  The
  same argument using \(z_2\) shifts \(a\) upward and proves
  non-local-nilpotence for \(z_2\).  Hence
  \(M_\hbar((\hbar))\) cannot be isomorphic to \(\mc P((\hbar))\) under
  the same Hamiltonian action.
\end{proof}

\begin{rmk}[Comparison with the heat-kernel propagator route]
  \label{rmk:no-heat-kernel-needed}
  The chain-level construction above avoids any heat-kernel
  propagator $P_{\epsilon,L}$ on $\R^2\times\C^2$.  In the
  Costello--Gwilliam BV propagator approach, the primitive
  indecomposable shadow is realized as a homotopy retraction onto the cohomology
  of the free open observable algebra, with the homotopy given by
  $G_L=\int_0^L K_t\,dt$ (the partial heat kernel).  Here the role of
  $G_L$ is played by the symmetric-algebra splitting, and the key
  analytic fact (positive jets are $Q$-exact,
  Lemma~\ref{lem:derivative-jets}) is needed only to
  identify the $E_1$-stalk of the brane factorization algebra with its
  $C^\infty$-smearing.  The two routes agree on the classical
  cocycle level by uniqueness of the cumulant projection on a
  symmetric algebra.

  The chain-level construction has the advantage of being independent
  of the analytic input of the heat kernel, and therefore independent
  of the Tate-completion convergence question recorded in
  Problem~\ref{prob:formal-factorization-center}, item (1).  It does
  not, however, by itself produce the renormalized BV homotopies
  required for the quantum (one-loop and higher) lift; that requires
  Costello renormalization of the brane-defect propagator, recorded
  separately in
  Problem~\ref{prob:formal-factorization-center}, items (1)--(2).
\end{rmk}

\begin{cor}[Primitive indecomposable shadow of the reduced cotangent comparison]
  \label{cor:unreduced-d-closed}
  Combining
  Theorem~\ref{thm:chain-level-primitive-projection} with the
  cochain-level closed-open derived center
  (Theorem~\ref{thm:open-closed-derived-center}) and the Hamiltonian
  factorization-algebra equivalence
  (Theorem~\ref{thm:open-closed-derived-center-factorization}), one
	  obtains a chain-level primitive indecomposable shadow of the reduced cotangent
  comparison:
  \begin{align*}
    \mathrm{Obs}^{\mathrm{cl}}_{\mathrm{closed},H}\big|_{L}
    &\xrightarrow{\;\Phi_{\mathrm{cl}}\;}
    \mathrm Z^{P_0}_{\mathrm{red,Ham}}\bigl(A_\partial^{\mathrm{pp}}\bigr) \\
    &\xrightarrow{\;\pi_{\mathrm{prim}}^{\mathrm{cot}*}\;}
    \text{primitive indecomposable shadow of the open boundary sector}.
  \end{align*}
  This corollary does not construct the unreduced cotangent
  factorization-center morphism. The missing data are the centrality
  homotopies for principal-part currents against arbitrary unreduced
  open observables, together with the Tate/weighted coefficient kernel
  and the brane-defect propagator extension recorded in
  Problem~\ref{prob:formal-factorization-center}.
\end{cor}

\begin{proof}
  $\Phi_{\mathrm{cl}}$ is constructed at the cochain level by
  Theorem~\ref{thm:open-closed-derived-center} on the reduced sector
  after de Rham-current contraction.  Theorem~\ref{thm:open-closed-derived-center-factorization}
  promotes this to a $P_0$ factorization algebra equivalence
  in the Hamiltonian sector.
	  Theorem~\ref{thm:chain-level-primitive-projection} gives the
	  primitive indecomposable \(P_0\)-shadow and one-$\psi$ label shadow. The
  composition therefore exists only at the indecomposable-shadow level.
  Since the one-$\psi$ projection does not intertwine the descendant
  raising action with the principal-part coadjoint lowering action, the
  asserted result is exactly the shadow statement and nothing stronger.
\end{proof}

\begin{rmk}[Status of the primitive indecomposable shadow]
  \label{rmk:T5-primitive-shadow-status}
	  The primitive indecomposable \(P_0\)-shadow map of
	  Theorem~\ref{thm:chain-level-primitive-projection} supplies the
	  local primitive indecomposable \(P_0\)-shadow component of
  Problem~\ref{prob:formal-factorization-center}, item~(3), and the
	  primitive indecomposable shadow component of
  Theorem~\ref{thm:main-local}. It does not by itself close
  Problem~\ref{prob:formal-factorization-center},
  because the Tate-completion of the Casimir
  ($\bigoplus_{d\geq 1}H_d^\vee\to$ a topological coefficient in
  which the Casimir converges) and the propagator extension are
  separate analytic obligations not addressed here. Nor does it supply
  the missing centrality homotopies for principal-part currents against
  arbitrary unreduced open observables.
  Proposition~\ref{prop:quantum-descendant-label-obstruction} further
  shows that the one-\(\psi\) label map is not compatible with the
  Moyal coadjoint action after quantization. The chain-level
  construction is therefore sufficient only for the classical primitive
  indecomposable shadow.
\end{rmk}

\begin{rmk}[On the role of the de Rham contraction]
  \label{rmk:role-of-deRham-contraction}
  The de Rham contraction Lemma~\ref{lem:derivative-jets} enters
  Construction~\ref{constr:chain-level-primitive-projection} only
  through the compactly supported density class on the brane line, not
  through any analytic propagator.  Specifically,
  \(H^\bullet(\Omega^\bullet_c(I))=H^\bullet_c(I;\C)\) is one-dimensional
  in degree \(1\) for an open interval \(I\), represented by a compactly
  supported density \(\lambda_I(t)\,dt\) with integral \(1\). The
	  degree-zero notation in the primitive shadow comes from the shifted
  BV/cotangent coordinate paired with this density; it is not ordinary
  point-evaluation in \(H_c^0(I)\).  The cosheaf structure is then
  transparent.
  No further analytic input is required; in particular no
  smoothing kernel, mollifier, or finite-scale BV gauge fix appears.
\end{rmk}
